\section{Related Work}
\label{sec:related_work}

The design of efficient multi-task architectures for resource-constrained edge computing intersects model merging, post-training quantization, dynamic routing, and TinyML systems.

\textbf{Model Merging and Parameter Fusion:} Parameter-space model merging combines specialized models without training from scratch using weight averaging \cite{weight_averaging}, permutation alignment \cite{git_rebasin}, Fisher approximations \cite{fisher_merging}, evolutionary search \cite{evolutionary_merging}, or task subspace edits \cite{concrete_subspaces}. For fine-tuned LoRA adapters \cite{lora}, Task Arithmetic \cite{task_arithmetic} adds task vectors, but suffers from sign conflicts. TIES-Merging \cite{ties_merging} resolves sign conflicts and prunes small weights, while DARES \cite{dares} uses dropout-based pruning. Other methods align features \cite{zipit} or operate in tangent space \cite{tangent_arithmetic}. However, weight-space merging is fundamentally \emph{static}: it cannot handle mixed-task batches at test-time or mitigate batch-size heterogeneity collapse \cite{regcalmerge}.

\textbf{Post-Training Quantization (PTQ):} PTQ compresses edge models by mapping weights and activations to low-bit integers (INT8/INT4) \cite{jacob2018integer, gholami2021survey}. Standard techniques mitigate outliers via channel scaling (SmoothQuant \cite{smoothquant}), second-order optimization (GPTQ \cite{gptq}), activation-aware scaling (AWQ \cite{awq}), dense-sparse decomposition (SqueezeLLM \cite{squeezellm}), weight clipping (OmniQuant \cite{omniquant}), or extreme 1-bit setups (BitNet \cite{bitnet}). However, direct PTQ of parameter-merged models collapses accuracy because merging disparate parameter manifolds creates extreme activation outliers that squash specialized features under coarse integer grids.

\textbf{Quantization-Aware Model Merging:} To bridge merging and quantization, Q-Merge \cite{qmerge} optimizes merging coefficients under quantization constraints using Straight-Through Estimators (STEs). While effective, Q-Merge requires expensive pseudo-gradient optimization over calibration data and suffers from severe cross-schema overfitting when deployed on target hardware that deviates from the training-time quantization operator.

\textbf{Dynamic Routing and Activation Blending:} Dynamic Mixtures of Experts (MoEs) route inputs to specialized sub-networks at test-time \cite{moe, switch_transformer} using conditional computation \cite{gshard}, balanced routing \cite{base_layers}, expert choice \cite{expert_choice}, or soft routing \cite{soft_moe, vmoe, coe}. For edge devices, SPS-ZCA \cite{sps_zca} and SABLE \cite{sable} route activations sample-wise using frozen adapters and pre-computed task centroids, bypassing weight-space interference and batch-size collapse. However, existing activation-blending methods require full precision (FP16/FP32). SA-QAB resolves this by introducing the first scale-aligned quantized activation blending framework.

\textbf{TinyML and Edge Systems:} Physical edge microcontrollers operate under severe memory constraints (e.g., $<256$KB SRAM) \cite{tinyml_survey}. Industry libraries like TensorFlow Lite Micro \cite{tflite_micro} provide integer execution kernels, while frameworks like MCUNet \cite{mcunet, mcunet_v2} and TinyEngine \cite{tinyengine} co-design networks and compilers to minimize memory. SA-QAB aligns with this ecosystem by employing diagonal-covariance GMMs and hard-argmax path pruning, fitting multi-task execution comfortably on microcontroller chips.
